\documentclass[10pt]{article}

% ---------- Document Identity (update these only) ----------
\newcommand{\DocStage}{}
\newcommand{\DocTitle}{Your Road to Tier One \textbf{SOF}}
\newcommand{\ProgramName}{182-Week Civilian-to-Selection Readiness Pipeline}
\newcommand{\ProgramSubtitle}{}

% ---------- Page ----------
\usepackage[legalpaper,left=0.5in,right=0.5in,top=0.6in,bottom=0.45in]{geometry}

% ---------- Typography ----------
\usepackage[T1]{fontenc}
\usepackage[utf8]{inputenc}
\usepackage{libertinus}
\usepackage{microtype}
\usepackage{parskip}
\usepackage{ragged2e}
\usepackage{textcomp}
\AtBeginDocument{\color{black}}
\AtBeginDocument{\pagecolor{white}}
\AtBeginDocument{\justifying}

% ---------- Landscape pages ----------
\usepackage{pdflscape}

% ---------- Color ----------
\usepackage[table]{xcolor}
\definecolor{StageBlue}{HTML}{1F4E79}
\definecolor{StageBlueDark}{HTML}{163A5A}
\definecolor{LightGray}{HTML}{F2F4F7}
\definecolor{MidGray}{HTML}{D9DEE5}
\definecolor{RuleGray}{HTML}{C7CED8}
\definecolor{HeaderInk}{HTML}{000000}
\definecolor{Ink}{HTML}{000000}

% ---------- Utility ----------
\usepackage{enumitem}
\sloppy
\setlength{\emergencystretch}{2em}

% ---------- Boxes / UI ----------
\usepackage[most]{tcolorbox}

\tcbset{
  charterTitle/.style={
    enhanced,
    colback=StageBlue!4,
    colframe=StageBlue,
    boxrule=1.25pt,
    arc=3pt,
    left=16pt,right=16pt,top=14pt,bottom=14pt,
    borderline north={3.2pt}{0pt}{StageBlue},
    borderline west={4pt}{0pt}{StageBlue},
    borderline south={1.25pt}{0pt}{StageBlue},
    drop shadow={black!25!white},
  },
  charterRule/.style={
    enhanced,
    colback=LightGray,
    colframe=RuleGray,
    boxrule=0.85pt,
    arc=3pt,
    left=12pt,right=12pt,top=10pt,bottom=10pt,
    borderline west={3pt}{0pt}{StageBlue},
  },
  charterBadge/.style={
    enhanced,
    colback=StageBlue,
    coltext=white,
    colframe=StageBlueDark,
    boxrule=0.6pt,
    arc=2pt,
    left=6pt,right=6pt,top=3pt,bottom=3pt,
  }
}
\newtcbox{\Badge}{nobeforeafter,charterBadge}

% ---------- Lists ----------
\setlist[itemize]{leftmargin=1.5em, itemsep=0.25em, topsep=0.25em}
\setlist[enumerate]{leftmargin=1.8em, itemsep=0.25em, topsep=0.25em}

% ---------- TOC ----------
\usepackage{tocloft}
\setcounter{tocdepth}{3}
\setlength{\cftbeforesecskip}{3pt}
\setlength{\cftbeforesubsecskip}{2pt}
\setlength{\cftbeforesubsubsecskip}{0pt}
\renewcommand{\cftsecfont}{\bfseries}
\renewcommand{\cftsecpagefont}{\bfseries}
\renewcommand{\cftsubsecfont}{\normalfont}
\renewcommand{\cftsubsecpagefont}{\normalfont}
\renewcommand{\cftsubsubsecfont}{\normalfont}
\renewcommand{\cftsubsubsecpagefont}{\normalfont}
\setlength{\cftsecindent}{0pt}
\setlength{\cftsubsecindent}{1.25em}
\setlength{\cftsubsubsecindent}{2.8em}
\setlength{\cftsecnumwidth}{2.1em}
\setlength{\cftsubsecnumwidth}{2.8em}
\setlength{\cftsubsubsecnumwidth}{3.6em}

% Remove the built-in TOC heading so only the custom heading is shown
\makeatletter
\renewcommand{\tableofcontents}{\@starttoc{toc}}
\makeatother

% ---------- Header/Footer: page number only ----------
\usepackage{fancyhdr}
\pagestyle{fancy}
\fancyhf{}
\renewcommand{\headrulewidth}{0pt}
\renewcommand{\footrulewidth}{0pt}
\fancyfoot[C]{\thepage}

% ---------- Section formatting ----------
\usepackage{titlesec}

\titleformat{\section}
  {\Large\bfseries\color{Ink}}
  {\thesection}{0.6em}{}
  [\vspace{0.25em}\textcolor{StageBlue}{\rule{\linewidth}{0.8pt}}]

\titleformat{\subsection}
  {\large\bfseries\color{Ink}}
  {\thesubsection}{0.6em}{}

\titleformat{\subsubsection}
  {\normalsize\bfseries\color{Ink}}
  {\thesubsubsection}{0.6em}{}

\titlespacing*{\section}{0pt}{0.95em}{0.35em}
\titlespacing*{\subsection}{0pt}{0.65em}{0.25em}
\titlespacing*{\subsubsection}{0pt}{0.55em}{0.2em}

% ---------- Table engine ----------
\usepackage{tabularray}
\UseTblrLibrary{booktabs}

% ---------- tabularray: GLOBAL table treatment ----------
\SetTblrInner{
  cells = {fg=black, bg=white, valign=t},
}

% ---------- Header helpers ----------
\newcommand{\Hdr}[1]{\shortstack[c]{\strut #1\strut}}

% ---------- Lines ----------
\newcommand{\TitleLine}{\textcolor{StageBlue}{\rule{0.98\linewidth}{0.95pt}}}
\newcommand{\SubTitleLine}{\textcolor{RuleGray}{\rule{0.98\linewidth}{0.65pt}}}

% ---------- Continued on next page ----------
\newcommand{\ContinuedOnNextPageInline}{%
  \par
  \vfill
  \noindent\textcolor{RuleGray}{\rule{\linewidth}{1pt}}\vspace{-2pt}
  \noindent\hfill\textit{Continued on next page}\par\vspace{-10pt}
  \noindent\textcolor{RuleGray}{\rule{\linewidth}{1pt}}
  \clearpage
}

% ---------- PDF hyperlinks ----------
\usepackage[hidelinks]{hyperref}
\hypersetup{
  colorlinks=false,
  pdfborder={0 0 0},
  linktoc=all,
  pdftitle={\DocTitle},
  pdfauthor={\ProgramName}
}

% =========================================================
\begin{document}

\section*{Stage 11.D — Nutritional Adjustments for Final Selection}
\phantomsection
\addcontentsline{toc}{section}{Stage 11.D — Nutritional Adjustments for Final Selection}

\subsection*{11.D.1 Purpose and Scope}
\phantomsection
\addcontentsline{toc}{subsection}{11.D.1 Purpose and Scope}

11.D defines how to apply the Stage 10 Nutrition Operating System in the final selection approach period so recovery, performance, and test validity remain stable during high-consequence phases and protected gate windows. The purpose is to reduce variance, prevent last-minute nutrition-driven validity contamination, and ensure that admissible performance evidence is not confounded by avoidable nutrition and hydration drift.

Exclusions:

\begin{itemize}
\item No meal plans, recipes, or grocery lists.
\item No new macro rules beyond Stage 10.
\item No supplement dosing, stacks, cycling, procurement guidance, or product recommendations.
\end{itemize}

\newpage

\subsection*{11.D.2 Late-Phase Nutrition Posture — Implementation Only}
\phantomsection
\addcontentsline{toc}{subsection}{11.D.2 Late-Phase Nutrition Posture — Implementation Only}

\begin{itemize}
\item Late phases (Phase 11–Phase 13) \textbf{MUST} apply Stage 10 macro posture with a stability-first bias:
\begin{itemize}
\item reduce volatility in calories and macronutrient distribution,
\item protect performance and recovery as consequence rises,
\item avoid aggressive deficit behavior when readiness and validity are primary objectives.
\end{itemize}
\item Trend-based decisions only:
\begin{itemize}
\item decisions \textbf{MUST} be based on weekly bodyweight average and weekly circumference suite trends plus recovery markers, not single-day scale movement.
\end{itemize}
\item Changes \textbf{MUST} comply with Stage 10.2 controls:
\begin{itemize}
\item calorie step size limited to ±150–250 kcal/day for structural adjustments,
\item hold period of 7–14 days after macro changes unless safety requires reversal,
\item structural calorie changes require ≥12/14 on-plan days,
\item minor tweaks require ≥6/7 on-plan days.
\end{itemize}
\item Deload/Test weeks \textbf{MUST} follow the Stage 10.2 binding rule:
\begin{itemize}
\item calories set to maintenance.
\end{itemize}
\item If recovery markers degrade or restrictive behaviors appear:
\begin{itemize}
\item maintenance reset is applied as defined in Stage 10.2 (7 days at maintenance),
\item exit the reset only when readiness stabilizes using Stage 2 recovery markers,
\item do not intensify restriction in response to late-phase anxiety.
\end{itemize}
\end{itemize}

\newpage

\subsection*{11.D.3 Test-Window Execution Rules — Stage 3 Alignment}
\phantomsection
\addcontentsline{toc}{subsection}{11.D.3 Test-Window Execution Rules — Stage 3 Alignment}

Operational rules for protected windows (implementation of Stage 10 and Stage 3; no new rules created):

Deload/Test weeks (protected windows)

\begin{itemize}
\item Calories \textbf{MUST} be set to maintenance per Stage 10.2.
\item Protein posture remains stable per Stage 10.A (do not reduce protein to “make room”).
\item Fat remains within the 20–30% floor band except for the test-day carb redistribution offset mechanism.
\item Test-day carb redistribution is allowed per Stage 10.2:
\begin{itemize}
\item +50–100 g carbs on test day(s),
\item offset equal calories reduced from fats,
\item no novelty foods/products.
\end{itemize}
\item No novelty posture applies:
\begin{itemize}
\item no new foods, no new electrolyte products/concentrations/timing practices, no new supplement categories, no caffeine pattern changes.
\end{itemize}
\end{itemize}

Gate proximity posture (final 7–10 days before major gate windows)

\begin{itemize}
\item No major diet structure change is permitted.
\item Caffeine pattern changes are prohibited within the Stage 10 novelty window:
\begin{itemize}
\item no caffeine pattern changes within 14 days of gate windows.
\end{itemize}
\item Any unavoidable logistical change \textbf{MUST} be logged with date, reason, and expected effect and treated as a validity confounder for conservative interpretation.
\end{itemize}

Reslot posture

\begin{itemize}
\item If nutrition disruption (GI distress, dehydration/heat illness risk, severe sleep collapse) contributes to invalid testing conditions, tests are reslotted to the next protected window per Stage 3.
\item No stacking make-ups; do not compress missed test attempts into unsafe density.
\end{itemize}

\begingroup
\scriptsize
\setlength{\tabcolsep}{2pt}
\renewcommand{\arraystretch}{1.12}

\begin{longtblr}{
width=\linewidth,
rowhead=1,
colspec = {
X[l,2.10]
X[l,2.10]
X[l,3.05]
X[l,2.85]
X[l,2.70]
},
hlines,
vlines,
cells = {hsep=6pt, vsep=6pt, halign=l, valign=t},
row{odd} = {bg=white},
row{even} = {bg=LightGray},
row{1} = {bg=StageBlue!15, fg=black, font=\bfseries, halign=l, valign=m},
}
\Hdr{Week Type} &
\Hdr{Goal} &
\Hdr{Allowed Changes} &
\Hdr{Not Allowed} &
\Hdr{What to Log} \
Normal training week (late phases) &
Stable performance support with low variance &
Minor distribution tweaks if ≥6/7 on-plan; structural changes only if ≥12/14 on-plan and not near gate window &
Single-day reactive cuts; multi-change weeks; novelty changes near gates &
Weekly \textbf{BW} avg + waist/circ suite; on-plan count; sleep/fatigue; any change date/hold/review \
Deload/Test week (Stage 3 protected) &
Validity protection and recovery &
Calories to maintenance; test-day carb redistribution (+50–100 g carbs offset by fats); protein stable; hydration stable per 10.3 &
New restrictive cuts; new foods/products; new supplement category; new electrolyte routine; caffeine pattern changes &
Maintenance posture confirmation; test-day carb offset Y/N; novelty Y/N; hydration notes \
7–10 days pre-gate window &
Controlled-variable stability &
\textbf{HOLD} posture; only safety-required corrections; maintain stable routine &
Any major macro shift; new food class; fiber overhaul; new supplement category; caffeine pattern change &
Gate proximity Y/N; novelty events; reason for any unavoidable change \
Post-test week (immediate) &
Resume stable execution without rebound extremes &
Resume prior posture only if recovery markers stable; maintenance reset if recovery collapsed &
Aggressive deficit “payback” &
Recovery markers; decision and hold period \
\end{longtblr}

\endgroup

\newpage

\subsection*{11.D.4 Travel and Access Instability Posture — Late Phase}
\phantomsection
\addcontentsline{toc}{subsection}{11.D.4 Travel and Access Instability Posture — Late Phase}

This section implements Stage 10.2 life-constraint posture in final phases with strict validity protection and logging continuity.

Rules (binding):

\begin{itemize}
\item Minimum viable execution is permitted:
\begin{itemize}
\item prioritize protein adequacy and calorie stability using simple, repeatable staples.
\end{itemize}
\item Portable baseline concept:
\begin{itemize}
\item use a small set of repeatable, predictable foods; reduce decision points.
\end{itemize}
\item Hydration continuity \textbf{MUST} follow Stage 10.3:
\begin{itemize}
\item baseline plus session top-ups,
\item electrolyte activation only when triggers apply,
\item no new electrolyte products or mixing practices in protected windows.
\end{itemize}
\item Log continuity \textbf{MUST} follow Stage 2.E:
\begin{itemize}
\item if fields cannot be recorded, record \textbf{UNKNOWN}/\textbf{MISSING}; no inference.
\end{itemize}
\item If travel disrupts a planned protected window:
\begin{itemize}
\item reslot the affected test to the next protected window; do not compress or stack.
\end{itemize}
\end{itemize}

\begingroup
\scriptsize
\setlength{\tabcolsep}{2pt}
\renewcommand{\arraystretch}{1.12}

\begin{longtblr}{
width=\linewidth,
rowhead=1,
colspec = {
X[l,2.35]
X[l,2.55]
X[l,2.55]
X[l,2.90]
X[l,2.65]
},
hlines,
vlines,
cells = {hsep=6pt, vsep=6pt, halign=l, valign=t},
row{odd} = {bg=white},
row{even} = {bg=LightGray},
row{1} = {bg=StageBlue!15, fg=black, font=\bfseries, halign=l, valign=m},
}
\Hdr{Constraint} &
\Hdr{Immediate Action} &
\Hdr{72-Hour Posture} &
\Hdr{7-Day Stabilization Posture} &
\Hdr{When to Escalate} \
Travel disruption (hotel/work) &
Activate portable baseline; maintain hydration baseline &
Hold macro changes; avoid novelty foods; log compliance daily &
Resume stable routine; if under-eating/fatigue rises, apply maintenance reset per Stage 10.2 &
Red-flag symptoms; persistent severe sleep collapse; repeated missed logs affecting admissibility \
Food access instability &
Simplify to staples; protect protein and calories &
Hold deficit escalation; focus on consistency over optimization &
Stabilize intake; structural changes deferred until ≥12/14 on-plan days achieved &
Restrictive behavior risk; persistent under-eating; illness flags \
Schedule volatility (late nights/early starts) &
Maintain stable meal timing where possible; avoid late boluses &
Hydration distribution rule; avoid caffeine pattern changes near gates &
If recovery markers degrade, maintenance reset 7 days &
Severe fatigue trend; repeated \textbf{MODIFIED}/\textbf{INVALID} sessions due to collapse \
Facility limitations (no prep tools) &
Operate within 7.13 minimal tooling posture &
Assembly posture; no meal plan; avoid new foods &
Re-establish stable routine as soon as tools return &
GI distress or dehydration symptoms persist \
\end{longtblr}

\endgroup

\newpage

\subsection*{11.D.5 GI Risk and Novelty Control — Validity Protection}
\phantomsection
\addcontentsline{toc}{subsection}{11.D.5 GI Risk and Novelty Control — Validity Protection}

Operational rules (implementation of Stage 10.4 GI ladder and Stage 3 no-novelty posture):

\begin{itemize}
\item Fiber and unfamiliar foods \textbf{MUST} remain stable near tests and gate windows:
\begin{itemize}
\item do not introduce new high-fiber loads, new sweeteners, or unusual volume spikes in the final 7–10 days before major gates.
\end{itemize}
\item Avoid novelty that commonly disrupts GI tolerance:
\begin{itemize}
\item new sugar alcohols,
\item novel “performance foods,”
\item unusually high-fat novelty meals,
\item sudden large increases in raw vegetables/legumes if not already tolerated.
\end{itemize}
\item Severe GI distress is a safety and validity disruptor:
\begin{itemize}
\item follow Stage 10.4 GI distress ladder posture,
\item reslot tests if symptoms compromise validity or safety,
\item do not “push through” gate attempts under active GI distress.
\end{itemize}
\end{itemize}

Final-phase Do Not Do list (GI/validity specific):

\begin{itemize}
\item Do not introduce new foods/products during Deload/Test weeks.
\item Do not change electrolyte product type or concentration during protected windows.
\item Do not use last-minute restrictive cuts to “feel lighter” for tests.
\item Do not change caffeine patterns within 14 days of gate windows.
\item Do not interpret acute GI distress as a reason to add multiple new changes; simplify and stabilize first.
\end{itemize}

\newpage

\subsection*{11.D.6 Optional Supplements Interface — Category Links Only}
\phantomsection
\addcontentsline{toc}{subsection}{11.D.6 Optional Supplements Interface — Category Links Only}

\begin{itemize}
\item Supplements are optional, category-level references only:
\begin{itemize}
\item 7.03 Foundational Supplements (optional),
\item 7.04 Additional Supplements / auxiliary session adjuncts (optional),
\item 7.05 Cognitive & Energy Enhancers (optional; do not change caffeine patterns near gates),
\item 7.06 Recovery & Sleep Aids (optional).
\end{itemize}
\item Dosing guidance and stack prescriptions are prohibited in Stage 11.D.
\item No novelty near gates:
\begin{itemize}
\item do not introduce new supplement categories, new products, or new timing practices inside Deload/Test weeks or within the final 7–10 days before major gates.
\end{itemize}
\item 7.02 Hormone Support & Adaptogens is education-only and never required.
\end{itemize}

Requires Stage 8 review flags (if encountered):

\begin{itemize}
\item Any upstream phase card or implementation artifact that treats supplements as required for deployability.
\item Any upstream artifact that implies 7.02 is required.
\end{itemize}

\newpage

\subsection*{11.D.7 Logging and Change Control — Late Phase}
\phantomsection
\addcontentsline{toc}{subsection}{11.D.7 Logging and Change Control — Late Phase}

Required documentation (binding; implements Stage 10.2 and Stage 2.E):

\begin{itemize}
\item All nutrition changes \textbf{MUST} be logged using Stage 10.2 change-log posture:
\begin{itemize}
\item date,
\item phase/week,
\item change made,
\item trigger(s),
\item intended outcome,
\item hold period,
\item review date,
\item owner role.
\end{itemize}
\item Gate proximity field \textbf{MUST} be included:
\begin{itemize}
\item within 14 days of gate windows Y/N,
\item if Y: confirm no novelty posture was applied and explain any unavoidable deviation.
\end{itemize}
\item Compliance tracking \textbf{MUST} be maintained:
\begin{itemize}
\item ≥6/7 on-plan days for minor tweaks,
\item ≥12/14 on-plan days for structural calorie changes.
\end{itemize}
\item Missing-data posture is binding:
\begin{itemize}
\item record \textbf{UNKNOWN} / \textbf{MISSING}; no inference; do not reconstruct by memory for gate decision use.
\end{itemize}
\item If logging completeness fails near gates:
\begin{itemize}
\item default posture is \textbf{HOLD} for decisions dependent on the missing evidence; reslot tests per Stage 3 if validity is compromised.
\end{itemize}
\end{itemize}

Requires Stage 8 review flags (if encountered):

\begin{itemize}
\item Any attempt to apply a nutrition rule that contradicts Stage 10 (e.g., planned refeeds/diet breaks, extreme weight-loss rates, water cuts).
\item Any attempt to alter deload/test week maintenance posture or no-novelty posture.
\end{itemize}

\newpage

\subsection*{11.D.8 Compliance Checklist — Pass/Fail}
\phantomsection
\addcontentsline{toc}{subsection}{11.D.8 Compliance Checklist — Pass/Fail}

\begin{itemize}
\item \textbf{PASS}/\textbf{FAIL}: No meal plans, recipes, or grocery lists included.
\item \textbf{PASS}/\textbf{FAIL}: No supplement dosing, stacks, cycling, procurement, or product recommendations included.
\item \textbf{PASS}/\textbf{FAIL}: No new macro rules introduced; Stage 10 is implemented only (10.A–10.D referenced by posture, not rewritten).
\item \textbf{PASS}/\textbf{FAIL}: Deload/Test weeks maintenance-calorie posture is stated and treated as binding; test-day carb redistribution posture is referenced as allowed where applicable.
\item \textbf{PASS}/\textbf{FAIL}: No novelty near gates posture is stated and enforced, including caffeine novelty window (no caffeine pattern changes within 14 days of gate windows).
\item \textbf{PASS}/\textbf{FAIL}: Logging and change-control posture is stated and aligned to Stage 10.2 and Stage 2.E missing-data doctrine (\textbf{UNKNOWN}/\textbf{MISSING}; no inference).
\item \textbf{PASS}/\textbf{FAIL}: Travel/access instability and GI risk posture included without redefining Stage 10; reslot posture per Stage 3 is explicit.
\item \textbf{PASS}/\textbf{FAIL}: Any mismatch is flagged "Requires Stage 8 review" in 11.D.6 and 11.D.7 rather than silently fixed.
\end{itemize}

\end{document}